%%%%%%%%%%%%%%%%%%%%%%%%%%%%%%%%%%%%%%%%%%%%%%%%%%%%%%%%%%%%%%%%%	
%	Nome:		cb_[].tex										%
%	Autor: 		Rafael Kümmel									%
%	Data:		19/09/20                						%
%	                              								%
%	Descrição: Arquivo principal que cria a apostila ciclo		%
%				básico                    						%
%                                                 			  	%	
%%%%%%%%%%%%%%%%%%%%%%%%%%%%%%%%%%%%%%%%%%%%%%%%%%%%%%%%%%%%%%%%%                            

\documentclass[a4paper, 11pt, twoside]{book}

% Tem que usar o pacote import antes, se não quebra

\usepackage{import}
\import{../../Universais/}{pacotes.sty}

%TODO: modificar com o número da apostila que está sendo criada!
\newcommand{\numapostila}{1}		% Número da apostila a ser criada

%TODO: Colocar a quantidade de capítulos de cada matéria aqui!
\newcommand{\capBio}{3}				% Quantidade de capítulos em Biologia 
\newcommand{\capFis}{3}				% Quantidade de capítulos em Fisica 
\newcommand{\capGeo}{3}				% Quantidade de capítulos em Geografia 
\newcommand{\capHis}{3}				% Quantidade de capítulos em Historia 
\newcommand{\capIng}{3}				% Quantidade de capítulos em Ingles 
\newcommand{\capMat}{3}				% Quantidade de capítulos em Matematica 
\newcommand{\capPor}{3}				% Quantidade de capítulos em Portugues 
\newcommand{\capQui}{3}				% Quantidade de capítulos em Quimica 

% ==========================================================================================

\begin{document}
	

	\frontmatter
	
	\import{../Formadores/}{capa_cb_0\numapostila.tex}
	\import{../../Universais/}{copyright.tex}	
	\import{../Formadores/}{homenagem\numapostila.tex}
	\import{../../Universais}{historiaCB.tex} 
	\import{../../Universais}{agradecimentos.tex} 
	\import{../Formadores}{autores\numapostila.tex}
	
	\pagecolor{fundo}
	\tableofcontents	
	%\thispagestyle{empty}
	
	%\listoffigures
	%\listoftables
	\mainmatter
	\pagecolor{fundo}
	% Aqui faz os capítulos da apostila, verificar a definicoes.sty	
	\begin{parte}{"Biologia é o estudo de coisas complicadas que têm a aparência de terem sido projetadas com um propósito"}{Richard Dawkins}{Biologia}
		\capitulos{\pathbiologia}{\capBio}
	\end{parte}
	
	\begin{parte}{“Dê-me um ponto de apoio fixo e eu faço a Terra mover-se.”}{Arquimedes}{Física}
		\capitulos{\pathfisica}{\capFis}
	\end{parte}
	
	\begin{parte}{"O geógrafo é, antes de tudo, um filósofo, e os filósofos são otimistas, porque diante deles está a infinidade."}{Milton Santos}{Geografia}
		\capitulos{\pathgeografia}{\capGeo}
	\end{parte}
	
	\begin{parte}{“A História Justifica o que se quer justificar. Ela não nos ensina nada a rigor, pois contém tudo e dá exemplos de tudo”}{Paul Valéry}{História}
		\capitulos{\pathhistoria}{\capHis}
	\end{parte}
	
	\begin{parte}{O ku tem trema}{Kümmel}{Inglês}
		\capitulos{\pathingles}{\capIng}
	\end{parte}
	
	\begin{parte}{"A matemática não mente. Mente quem faz mau uso dela."}{Albert Einstein}{Matemática}
		\capitulos{\pathmatematica}{\capMat}
	\end{parte}
	
	\begin{parte}{"Eu não escrevo em português. Escrevo eu mesmo."}{Fernando pessoa}{Português}
		\capitulos{\pathportugues}{\capPor}
	\end{parte}
	
	\begin{parte}{"A química é necessariamente uma ciência experimental: as conclusões são extraídas de dados e seus princípios são apoiados pela evidência dos fatos."}{Michael Faraday}{Química}
		\capitulos{\pathquimica}{\capQui}
	\end{parte}
	
	\backmatter
	
	\import{../../Universais}{nomenclatura.tex} 
	
	\printindex
	
	\nocite{*}
	\bibliographystyle{ieeetr}
	\bibliography{./Bibliografia/ref_Fis,%
		./Bibliografia/ref_Qui,%
		./Bibliografia/ref_Mat,%
		./Bibliografia/ref_Ing,%
		./Bibliografia/ref_Por,%
		./Bibliografia/ref_His,%
		./Bibliografia/ref_Geo,%
		./Bibliografia/ref_Bio}
	
\end{document}